\scanpage{220}\mspage{165}

% MS 1–2
déduite par ``transposition'' ($\mathrm{Hom}(-,b)$) de la flèche

% MS 3
\[
 p_{S'!}(a_{S'})\longrightarrow p_{S!}(a_S)
\]

% MS 4–6
induite par $f$ -- donc il faut prouver que celle-ci est un iso (c'est la forme que nous avions utilisée tantôt). Or cette flèche s'identifie à

% MS 7
\[
 \unsure{s|S}'\times a\longrightarrow s\times a
\]

% MS 8
induite par

% MS 9–10
\diagram{\begin{tikzcd}[ampersand replacement=\&,column sep=small,row sep=0.2em]
 \mathrm{hot}_{\unsure{W|w}}(f): \& s'\arrow[r,no head] \& s\\
 \& \shortparallel \& \shortparallel\\
 \& \mathrm{hot}_W(S') \& \mathrm{hot}_W(S)
\end{tikzcd}}

% MS 10
,

% MS 11–14
et dire que c'est iso quel que soit $a$, dans $\mathrm{Hot}_W$, revient à dire que $s'\longrightarrow s$ l'est (puisque la multiplication dans $\mathrm{Hot}_W$ a un élément unité \ldots)

% MS 15–22
Maintenant je me rappelle que pour avoir l'axiome Der 2 des dérivateurs, je devais également supposer $W$ \uline{fortement} \unsure{saturée|saturé} (i.e. égale au localiseur défini par le foncteur canonique $\mathrm{Cat}\longrightarrow(\mathrm{Cat})W^{-1}\overset{\mathrm{def}}{=}\mathrm{Hot}_W$). (\unsure{C'est} ce que je supposais d'emblée, en partant d'un $(M,W)$ \inserted{général}, pour définir un prédérivateur, pour avoir l'axiome Der 2.)

% MS 23–25
En appliquant le dual du th. de I, p.~\unsure{56|50}, on trouve donc le th. de chgt. de base pour $f_!$, pour chgt de base $p:S'\longrightarrow S$ \inserted{$W$-}\uline{propre}.

% archive sheet number: 219
